% sections/05_results.tex: model comparison, the final model and its selection-aware estimate, robustness,
% calibration, several attempts, failures, natural use, learning curve. Details: Appendix F (app:final).
\section{Results}\label{sec:results}

\subsection{Model comparison}\label{sec:results:comparison}
Figure~\ref{fig:FINAL_01_master_forest} shows the \R{final/n_candidates} candidates as their lanes deliver them.
Most of them beat the reference floor, yaw8 (\Rci{reference/accuracy}): \R{final/paired_ref_better} of
\R{final/paired_ref_m} candidates are more accurate after Holm adjustment. The differences among the best
candidates are small: the candidate with the lowest log loss and the others that paired comparisons cannot separate
from it form \R{final/concl_n_leaders} leaders (\R{final/paired_best_sep} of the \R{final/paired_best_m} other
candidates are separated from it, Holm over these \R{final/paired_best_m} comparisons). Of the candidates, \R{final/n_exploratory_candidates} were
chosen on the outer folds in their own lane, so the forest is exploratory. Appendices~\ref{app:m1}
to~\ref{app:m3} give each family's own results and decisions.

\begin{figure}[htbp]
\centering
\genfig{final/FINAL_01_master_forest}
\caption{The \R{final/n_candidates} candidates of the cross-lane selection as their lanes deliver them (LOSO
accuracy and log loss), coloured by lane. Open markers: chosen on the outer folds in their lane (exploratory);
diamonds: log loss after the calibration remedy chosen inside each fold; dashed: the reference floor. The forest is
exploratory; the nested estimate (Table~\ref{tab:final_nested}) is the honest number.}
\label{fig:FINAL_01_master_forest}
\end{figure}

\subsection{The final model and its selection-aware estimate}\label{sec:results:final}
The final model is the body-frame 1D CNN: a one-dimensional convolutional network on 12 channels in the gravity
frame and the body frame (Appendix~\ref{app:m3}). It has the lowest log loss in production form. The headline is the
nested, selection-aware estimate of the cross-lane selection, which reruns the rule inside every outer fold and
treats the \R{final/n_candidates} candidates as fixed (Table~\ref{tab:final_nested},
Figure~\ref{fig:FINAL_03_nested_selection}): accuracy \Rci{final/nested_accuracy}, AUC \Rcib{final/nested_auc} and
log loss \Rcib{final/nested_log_loss}. Inside the outer folds the selection chose the CNN in \R{final/sel_m3_deep}
of \R{final/sel_folds} folds. It does not repeat the lanes' own outer-fold choices: lane m3's own nested estimates,
which repeat its configuration choice in every fold, are \R{m3/nested_deep_nested_inv_vs_cnn_acc} among its
grip-invariant configurations and \R{m3/nested_deep_nested_all_vs_cnn_acc} among all of them (both post hoc;
Appendix~\ref{app:m3:selection}).

The shipped model's own LOSO numbers are exploratory: it was the best of \R{final/final_configs} configurations in
its lane, then chosen among \R{final/n_candidates} candidates. They are accuracy \Rci{final/final_accuracy}, AUC
\Rcib{final/final_auc} and log loss \Rcib{final/final_log_loss}, an optimism of \Rcib{final/optimism_accuracy}
points of accuracy over the nested estimate. Its body-frame input was also designed after outer-fold results had
been seen, an optimism that no nested estimate removes (Appendix~\ref{app:m3:selection}). The rest of this section
describes the shipped model, and its numbers share this exploratory status. The CNN is
\Rcib{final/final_vs_ref_acc}~points more accurate than the reference floor (Holm $p$~\R{final/final_vs_ref_holm_acc} over the
\R{final/paired_ref_m} comparisons with the reference).

\gentab{final/final_nested}

\begin{figure}[htbp]
\centering
\genfig{final/FINAL_03_nested_selection}
\caption{Nested cross-lane selection. Left: how often each candidate was chosen, by the lowest inner log loss on
the training persons, over the 20 outer folds. Right: the selection-aware estimate against the best candidate
picked on the outer folds. The naive pick is optimistic by \Rcib{final/optimism_accuracy}~points of accuracy.}
\label{fig:FINAL_03_nested_selection}
\end{figure}

The input and the model class were compared separately. On the same body-frame channels, the CNN is
\Rcib{final/ic_cnn_vs_m2_same12_acc}~points more accurate than lane m2's engineered features with logistic
regression (Holm $p$ = \R{final/ic_cnn_vs_m2_same12_holm_acc} over six input-against-class comparisons), and on the gravity frame alone it does not beat
lane m1's linear models (Appendices~\ref{app:final:input} and~\ref{app:m3:engineered}). With all 21 channels, lane
m2's logistic regression is not detectably less accurate than the CNN (\Rcib{final/ic_m2_all21_vs_cnn_acc}~points,
Holm $p$ = \R{final/ic_m2_all21_vs_cnn_holm_acc}) but has a higher log loss (\Rcib{final/ic_m2_all21_vs_cnn_ll},
Holm $p$ = \R{final/ic_m2_all21_vs_cnn_holm_ll}).

\subsection{Robustness to an unseen grip}\label{sec:results:robustness}
The robustness tests decided among the leaders (Figure~\ref{fig:FINAL_06_robustness}). Trained on upright grips
only and tested on the sideways grips, the leaders that read device-frame channels collapse: lane m2's logistic
regression changes by \Rb{final/rob_m2_lr_grip_d}~points, its stack by \Rb{final/rob_m2_stack_grip_d}, the
preprocessing panel's gradient-boosted trees by \Rb{final/rob_p1_gbm_grip_d} and lane m2's LightGBM by
\Rb{final/rob_m2_lgbm_grip_d} (each Holm $p \le$ \R{final/rob_m2_lgbm_grip_holm} over the \R{final/fig_grip_tested}
grip tests). The body-frame CNN changes by \Rcib{final/final_grip_d}~points (unadjusted $p$ =
\R{final/final_grip_p}; Holm $p$ = \R{final/rob_m3_deep_grip_holm} over the \R{final/fig_grip_tested} grip tests)
and reaches \R{final/rob_m3_deep_grip} on the unseen grips. It is one of the two leaders without a Holm-significant
drop; the other, an ensemble, changes by \Rb{final/rob_ens_mean3_grip_d}~points (Holm $p$ =
\R{final/rob_ens_mean3_grip_holm}). Of these two, the rule took the more portable CNN. The documented fallback is
lane m1's PHYS6 logistic: \Rb{final/simple_vs_best_acc}~points less accurate (Holm $p$ =
\R{final/simple_vs_best_holm_acc} over the comparisons with the best candidate), but it changes by only
\Rb{final/simple_grip_d}~points on the unseen grips (\R{final/simple_grip_acc}; Appendix~\ref{app:final:fallback}).

A phone held with the screen facing away defeats every model that reads the screen-relative frame: on synthesised
screen-away twins of the recorded trials the CNN falls to \R{final/final_flip_raw}. A model-free check of the rest
pose before the movement, exported with the model, detects \R{final/final_check_fires_flip} of the twins; with the
check and an un-flip the CNN recovers \R{final/final_flip_fixed}. The check also fires on
\R{final/final_check_fires_rec} of the recorded trials, which costs \Rcib{final/final_flip_cost}~points of accuracy
on them (Appendix~\ref{app:final:screen}).

\begin{figure}[htbp]
\centering
\genfig{final/FINAL_06_robustness}
\caption{Robustness of every candidate (bold: the final model). Left: accuracy on the sideways grips when trained
on every grip (dot) or on upright grips only (square). Middle: accuracy as recorded and on synthesised screen-away
twins, without and with the rest-pose check. Right: accuracy change when trained on the other platform only and
without suspect label blocks. Device-frame models fail an unseen grip, which removes them from the selection.}
\label{fig:FINAL_06_robustness}
\end{figure}

\subsection{Calibration}\label{sec:results:calibration}
The final model ends with a Platt calibrator fitted inside each fold (lane m3's choice, Appendix~\ref{app:m3}). With
it, its probabilities are consistent with calibration (Figure~\ref{fig:FINAL_08_calibration}): ECE
\R{final/final_cal_ece}, inside its perfect-calibration null band (up to \R{final/final_cal_ece_null_hi}); slope
\Rcib{final/final_cal_slope}; intercept \Rcib{final/final_cal_intercept}. Four further remedies fitted inside the
folds on top of it (Platt, temperature, isotonic and beta calibration) each raise the log loss slightly, none
detectably (Holm $p$ = \R{final/cal_platt_holm} for each, over the four remedies), so none is added. No difference between the platforms is detectable
(slope difference $p$ = \R{final/cal_platdiff_slope_p}, calibration in the large $p$ = \R{final/cal_platdiff_citl_p},
unadjusted), but with \R{final/cal_ios_persons} iOS persons this comparison has little power, and per-platform
intercepts were not saved (Appendix~\ref{app:final:calibration}).

\begin{figure}[htbp]
\centering
\genfig{final/FINAL_08_calibration}
\caption{Reliability of the final model, overall and per platform (equal-mass bins, person-bootstrap 95\,\% bands;
its in-fold Platt calibrator, no further remedy). ECE \R{final/final_cal_ece} (null band up to \R{final/final_cal_ece_null_hi}), slope
\Rcib{final/final_cal_slope}, intercept \Rcib{final/final_cal_intercept}: the percentage shown is consistent with
calibration.}
\label{fig:FINAL_08_calibration}
\end{figure}

\subsection{Several attempts}\label{sec:results:attempts}
Combining attempts helps (Figure~\ref{fig:FINAL_10_attempts}). With $k = 3$ attempts of one person and hand, the
mean \pright{} of the evaluation harness's $k$-curve reaches \Rci{final/final_kthree}. The aggregates of
Figure~\ref{fig:FINAL_10_attempts} give nearly the same answers (each is right in \R{final/agg_k3_sum_tau_acc}) and
differ in the percentage they show. The displayed aggregate of equation~\eqref{eq:aggregate}, with $\tau_3$ fitted
in-fold (median \R{final/agg_tau3}), is neither over- nor under-confident on average: its mean confidence minus its
accuracy is \Rcib{final/final_kthree_over}. This is calibration in the large only; its ECE with a null band, slope
and intercept were not computed. The mean \pright{} is under-confident (\Rcib{final/agg_k3_mean_p_over}) and has a
higher log loss (difference \R{final/agg_dll_k3} in favour of the scaled sum, Holm $p$ =
\R{final/agg_dll_k3_holm} over the paired comparisons of the aggregates; Appendix~\ref{app:final:attempts}).

\begin{figure}[htbp]
\centering
\genfig{final/FINAL_10_attempts}
\caption{Several attempts of one person and hand (\R{final/agg_persons} persons with at least 5 trials of each
hand). Left: over-confidence (mean confidence minus accuracy) of four aggregates (above 0: over-confident); right:
their log loss. The scaled sum of logits is neither over- nor under-confident on average; the mean \pright{} is
under-confident and is not shown to the user.}
\label{fig:FINAL_10_attempts}
\end{figure}

\subsection{Who and where it fails}\label{sec:results:failures}
Figure~\ref{fig:FINAL_11_persons} breaks the accuracy down. It is \R{final/brk_android_acc}
[\Rb{final/brk_android_lo}, \Rb{final/brk_android_hi}] on Android and \R{final/brk_ios_acc}
[\Rb{final/brk_ios_lo}, \Rb{final/brk_ios_hi}] on iOS; the platforms were not compared by a test. The accuracy is
below \R{final/person_threshold} for \R{final/brk_below80} persons, the lowest being \R{final/brk_low_person}; the
two left-handed persons score \R{final/brk_left_p11} and \R{final/brk_left_p17} (two persons: no CI). The weakest
grip-by-speed cell still reaches \Rci{final/brk_worst_cell}; the weakest movement variant reaches
\R{final/brk_worst_variant} [\Rb{final/brk_worst_variant_lo}, \Rb{final/brk_worst_variant_hi}] on
\R{final/brk_worst_variant_n} trials of \R{final/brk_worst_variant_p} persons. Of the \R{final/brk_err} errors,
\R{final/brk_conf_err} are confident errors (a probability below 0.2 for the true hand), and they cluster in a few
persons (\R{final/brk_conf_p09} in P09, \R{final/brk_conf_p20} in P20; Appendix~\ref{app:final:strata}).

\begin{figure}[htbp]
\centering
\genfig{final/FINAL_11_persons}
\caption{The final model by person (left; sorted; colour: platform; L: the two left-handed persons) and by grip
$\times$ speed cell (right; CI only for cells with at least 5 persons, open markers otherwise).}
\label{fig:FINAL_11_persons}
\end{figure}

\subsection{Natural use}\label{sec:results:natural}
On the free-choice trials the final model's right-hand recall is \R{final/final_nat_recall}
(\R{final/final_nat_credit} of \R{final/final_nat_n}), against \R{final/nat_m3_deep_exp} correct trials expected
from the same persons' instructed trials: no detectable shortfall (Holm $p$ = \R{final/nat_m3_deep_holm} over the
\R{final/n_candidates} candidates). All free-choice
trials are right-hand trials so far, so a constant right answer would score 100\,\% and this recall says nothing
about the left hand. The mean probability of the reported hand is lower on free choice than on the same
persons' instructed right-hand trials (\R{final/cal_nat_meanp} against \R{final/cal_nat_primary}), but the
difference is not detectable ($p$ = \R{final/cal_nat_p}, unadjusted; Appendix~\ref{app:final:natural}).

\subsection{More persons?}\label{sec:results:learning}
The learning curve flattens (Figure~\ref{fig:FINAL_12_learning_curve}). With \R{final/lc_m3_deep_nall} training
persons the CNN reaches \R{final/lc_m3_deep_accall}. An inverse power law fitted to random subsets of training
persons extrapolates to \Rci{final/lc_m3_deep_at40} at 40 persons, a gain along the fitted curve of
\Rcib{final/lc_m3_deep_gain40}~points from \R{final/lc_m3_deep_nall} persons. The band resamples the repetitions at
fixed persons, so it understates the uncertainty, and the curve is fitted to few sizes. The extrapolation suggests
that more persons of the same kind would add little average accuracy; new persons are needed above all for the
groups the data lack: left-handers and left-hand free choice.

\begin{figure}[htbp]
\centering
\genfig{final/FINAL_12_learning_curve}
\caption{Learning curves: training persons drawn at random inside every outer fold (dots: repetitions; line: their
mean; dashed: an inverse power law with a 95\,\% bootstrap band over the repetitions, which understates the
uncertainty). The extrapolation to 40 persons gives \R{final/lc_m3_deep_at40} for the CNN; by this extrapolation,
more persons of the same kind would add little average accuracy.}
\label{fig:FINAL_12_learning_curve}
\end{figure}
